\subsection{微调概述}

\noindent 预训练完成后，大型语言模型（LLM）可被直接用于解决各类 NLP 任务。其工作方式本质上是\textbf{文本补全}：给定一段上下文，模型生成最可能的后续文本。形式化地，设上下文为 token 序列 $\mathbf{x}=x_0...x_m$，待生成序列为 $\mathbf{y}=y_1...y_n$，则推理过程可定义为
\begin{eqnarray}
\hat{\mathbf{y}} & = & \argmax_{\mathbf{y}} \log \Pr(\mathbf{y} | \mathbf{x}) \nonumber \\
                 & = & \argmax_{\mathbf{y}} \sum_{i=1}^{n} \log \Pr(y_i|x_0,...,x_m,y_1,...,y_{i-1}) \label{eq:inf}
\end{eqnarray}
这与标准语言模型的对数概率形式完全一致，区别仅在于预测从位置 $m+1$ 开始。因此，只需将任务包装成合适的\textbf{提示模板}，就能让同一个 LLM 处理多种问题。例如，判断句子语法是否正确，可构造如下输入：

\vspace{0.5em}
\begin{tcolorbox}[frame empty]
\hspace{4em} John seems happy today.

\hspace{4em} 问题：这个句子语法正确吗？

\hspace{4em} 答案：\underline{\hspace{2em}}
\end{tcolorbox}
\vspace{0.5em}

模型在上下文后生成 ``Yes'' 即给出了判断。同样地，机器翻译可以使用 ``Translate this sentence from English into Chinese:'' 这样的指令模板，甚至可以写成类似代码的配置块。这种统一范式的强大之处在于，无需为每个任务设计专门的模型结构。

然而，让 LLM 稳定地“听懂”指令并非天然具备的能力。一种直接的思路是将大量（指令，响应）样本混入预训练数据，但这需要海量标注，且从头训练的计算代价极高。因此，目前主流且高效的做法是\textbf{指令微调}（instruction fine-tuning）：在预训练模型的基础上，用一套规模小得多但质量高的指令数据集 $\mathcal{D}_{\text{tune}}$ 对参数进行轻微调整，从而激活模型在预训练中已习得的通用语言知识，使其泛化到各种指令驱动的任务上。

\textbf{指令微调数据}由输入-输出对构成。输入部分融合了任务指令、可选的系统提示以及用户提供的内容，输出部分则为对应的标准回答。以下给出了几个典型的微调样本，它们覆盖了分类、生成、推理和代码等不同任务，并且同一任务（如语法判别）可以对应风格迥异的指令，从而增加数据多样性：

\vspace{0.5em}
\begin{tcolorbox}[frame empty]
\begin{center}
\parbox{0.9\textwidth}{
{\color{blue} 阅读以下句子，并回答其语法是否正确。}

{\color{blue} LLMs are powerful models but are expensive to build.}

\vspace{0.3em}
\underline{是}
}
\end{center}
\end{tcolorbox}

\vspace{0.5em}
\begin{tcolorbox}[frame empty]
\begin{center}
\parbox{0.9\textwidth}{
{\color{blue} 写一个关于三个角色的短篇故事：一个叫 Andy 的男孩、他的狗和他的朋友 Jane。}

\vspace{0.3em}
\uline{在繁华城市的中心，Andy、他忠实的狗 Rocket 和他的朋友 Jane 在暮色中发现了一个隐藏的花园……}
}
\end{center}
\end{tcolorbox}

\vspace{0.5em}
\begin{tcolorbox}[frame empty]
\begin{center}
\parbox{0.9\textwidth}{
{\color{blue} 如果你买 5 个苹果，每个苹果 1.20 美元，总共花多少钱？}

\vspace{0.3em}
\underline{6.00 美元}
}
\end{center}
\end{tcolorbox}

\vspace{0.5em}
\begin{tcolorbox}[frame empty]
\begin{center}
\parbox{0.9\textwidth}{
{\color{blue} 写一个 Python 程序来计算以下数字的平方和：1、2、10、-9、78。}

\vspace{0.3em}
\underline{\texttt{numbers = [1,2,10,-9,78]}}

\underline{\texttt{sum\_of\_squares = sum(x**2 for x in numbers)}}

\underline{\texttt{print(sum\_of\_squares)}}
}
\end{center}
\end{tcolorbox}
\vspace{0.5em}

这些示例中，输入部分（蓝色）在计算损失时会被屏蔽，模型仅学习预测下划线标注的输出部分。研究指出，扩大微调任务的数量和指令的多样性可以显著提升模型性能 \cite{chung-etal:2022scaling}。但值得强调的是，微调数据的体量远小于预训练语料——通常数万至数十万条高质量样本即已足够 \cite{zhou-etal:2023lima}，而预训练则需要数十亿甚至数万亿 token，两者在计算开销上差了几个数量级。

更为关键的是，经过指令微调后，模型往往表现出\textbf{零样本泛化}（\mindex{zero-shot learning}）能力：它不仅能完成训练时见过的指令类型，还能应对从未被显式训练过的全新任务描述 \cite{sanh-etal:2022multitask,wei-etal:2022finetuned}。这种能力正是生成式 LLM 与早期 BERT 等“针对每个下游任务独立微调”范式的分水岭。背后的原理通常被理解为，预训练已经储备了充足的语言理解和推理基础，指令微调只是引入了一种监督信号，将潜藏的指令遵循与任务解决能力“唤醒”。

\noindent \textbf{微调的数学形式}与预训练高度相似，但目标函数聚焦于输出段。设预训练得到的参数为 $\hat{\theta}$，微调数据集为 $\mathcal{D}_{\text{tune}}$，则最优参数 $\tilde{\theta}$ 为：
\begin{eqnarray}
\tilde{\theta} = \argmin_{\hat{\theta}^+} \sum_{\mathrm{sample} \in \mathcal{D}_{\text{tune}}} \mathcal{L}_{\hat{\theta}^+}(\mathrm{sample}) \label{eq:llm-fine-tuning-mle}
\end{eqnarray}
对每个样本，我们将其切分为输入段 $\mathbf{x}_{\text{sample}}$ 和输出段 $\mathbf{y}_{\text{sample}}$，损失函数只作用于后者：
\begin{eqnarray}
\mathcal{L}_{\hat{\theta}^+}(\mathrm{sample}) = -\log \Pr_{\hat{\theta}^+}(\mathbf{y}_{\text{sample}}|\mathbf{x}_{\text{sample}})
\end{eqnarray}
在具体实现中，前向传播照常对整个拼接序列进行计算，但反向传播时，误差梯度仅流经 $\mathbf{y}_{\text{sample}}$ 对应的位置。考虑这样一个序列：
\begin{equation}
    \underbrace{\textrm{{\color{blue} $\langle s \rangle$ 计算这个数的平方。}}\ \textrm{{\color{blue} $2$。}}}_{\text{输入}}\ \ \underbrace{\textrm{\underline{结果是 $4$。}}}_{\text{输出}}\nonumber
    \end{equation}
    损失仅基于 “结果是 $4$。” 这几个 token 计算并回传，输入部分不产生任何梯度。这种方式确保模型专注于学习如何根据指令生成正确的响应，而非重复记忆输入内容。
指令微调在实际操作中仍需要一定的工程调参——学习率、批大小、训练步数等超参数需要仔细摸索。不过，由于其数据规模和计算量远不及预训练，多次实验的整体成本是完全可控的。正是这种低成本、高效率的特性，使得微调成为适配 LLM 的通用技术栈核心：除了指令微调，它还被广泛用于将基座模型转化为对话助手、适配超长上下文、进行领域知识注入等。后续章节将深入探讨更高效的微调算法以及在更多场景下的应用。
